\documentclass[10pt, letterpaper]{article}

% ── Paquetes ──────────────────────────────────────────────────────────────────
\usepackage[T1]{fontenc}
\usepackage[utf8]{inputenc}
\usepackage[spanish]{babel}
\usepackage{lmodern}
\usepackage[left=1.8cm, right=1.8cm, top=1.5cm, bottom=1.5cm]{geometry}
\usepackage{titlesec}
\usepackage{enumitem}
\usepackage{hyperref}
\usepackage{xcolor}
\usepackage{tabularx}
\usepackage{array}
\usepackage{parskip}

% ── Colores ───────────────────────────────────────────────────────────────────
\definecolor{azul}{RGB}{20, 60, 120}
\definecolor{gris}{RGB}{100, 100, 100}
\definecolor{linea}{RGB}{180, 200, 230}

% ── Hipervínculos ─────────────────────────────────────────────────────────────
\hypersetup{colorlinks=true, urlcolor=azul, linkcolor=azul}

% ── Secciones ─────────────────────────────────────────────────────────────────
\titleformat{\section}
  {\large\bfseries\color{azul}}
  {}
  {0em}
  {}
  [\color{linea}\hrule]

\titlespacing{\section}{0pt}{8pt}{4pt}

% ── Listas compactas ──────────────────────────────────────────────────────────
\setlist[itemize]{noitemsep, topsep=2pt, partopsep=0pt,
                  leftmargin=1.2em, label={\small\textcolor{azul}{$\bullet$}}}

% ─────────────────────────────────────────────────────────────────────────────
\begin{document}
\pagestyle{empty}

% ═══════════════════════════════════════════════════════════════════════════════
%  ENCABEZADO
% ═══════════════════════════════════════════════════════════════════════════════
\begin{center}
  {\LARGE\bfseries\color{azul} Jair Molina Arce}\\[4pt]
  {\large\color{gris} Ingeniero Mecánico}\\[6pt]
  \small
  \textcolor{gris}{Ciudad de México, México}~\textbf{·}~
  \href{mailto:ingjairmolina@gmail.com}{ingjairmolina@gmail.com}~\textbf{·}~
  5652646108\\[2pt]
  CVU CONACYT: 1340773~\textbf{·}~ORCID: 0009-0009-6732-8100
\end{center}

\vspace{6pt}

% ═══════════════════════════════════════════════════════════════════════════════
%  PERFIL PROFESIONAL
% ═══════════════════════════════════════════════════════════════════════════════
\section{Perfil Profesional}

Ingeniero Mecánico egresado del Instituto Politécnico Nacional con sólida formación en
materiales, procesos de manufactura, instrumentación y análisis de sistemas.
Capacidad analítica para evaluar cotizaciones, comparar especificaciones técnicas de
materiales e insumos, y coordinar con proveedores para garantizar la calidad y oportunidad
de los suministros. Orientado al trabajo bajo presión, con disponibilidad de tiempo
completo y apertura a cambio de residencia en territorio nacional.

% ═══════════════════════════════════════════════════════════════════════════════
%  EDUCACIÓN
% ═══════════════════════════════════════════════════════════════════════════════
\section{Educación}

\noindent
\begin{tabularx}{\linewidth}{@{}X r@{}}
  \textbf{Maestría en Tecnologías Avanzadas} \textcolor{gris}{(en curso)} &
    \textcolor{gris}{2024 – Presente} \\
  \textit{Instituto Politécnico Nacional (IPN)} & \\[4pt]

  \textbf{Ingeniería Mecánica} &
    \textcolor{gris}{2017 – 2023} \\
  \textit{Instituto Politécnico Nacional (IPN)} & \\[4pt]

  \textbf{Bachillerato Técnico en Informática} &
    \textcolor{gris}{2014 – 2017} \\
  \textit{Colegio de Bachilleres Plantel 1 ``El Rosario''} & \\
\end{tabularx}

% ═══════════════════════════════════════════════════════════════════════════════
%  EXPERIENCIA PROFESIONAL
% ═══════════════════════════════════════════════════════════════════════════════
\section{Experiencia Profesional}

% --- Cargo 1 ---
\noindent
\textbf{Ingeniero de Sistemas e Instrumentación} \hfill
\textcolor{gris}{2022 – 2024}\\
\textit{Instituto Politécnico Nacional (IPN)} — Ciudad de México
\begin{itemize}
  \item Selección, adquisición y gestión técnica de sensores, actuadores e insumos
        (presión, temperatura, flujo) para banco de pruebas de propulsión.
  \item Evaluación de especificaciones técnicas de componentes electrónicos y mecánicos;
        comparativa de proveedores nacionales para reducir costos y garantizar calidad.
  \item Diseño e integración de sistemas de adquisición de datos (DAQ) con
        microcontroladores ESP32 y STM32; coordinación de logística de entrega de equipo.
  \item Análisis de datos operativos con Python, pandas y MATLAB/Simulink para validación
        de desempeño de insumos adquiridos.
\end{itemize}

\vspace{4pt}
% --- Cargo 2 ---
\noindent
\textbf{Desarrollador Full-Stack — Plataforma de Monitoreo Industrial} \hfill
\textcolor{gris}{2023 – 2024}\\
\textit{Proyecto Independiente / IPN} — Ciudad de México
\begin{itemize}
  \item Coordinación con fabricantes y distribuidores para la integración de dispositivos
        IoT en plataforma de monitoreo; levantamiento de requerimientos técnicos.
  \item Gestión de órdenes de compra de equipos de cómputo y electrónica; verificación de
        cumplimiento de especificaciones antes de aceptar recepción.
  \item Despliegue de infraestructura tecnológica (Docker, servidores, bases de datos)
        incluyendo cotización y contratación de servicios en la nube.
\end{itemize}

\vspace{4pt}
% --- Cargo 3 ---
\noindent
\textbf{Docente — Tecnologías Disruptivas (VR/AR)} \hfill
\textcolor{gris}{2024}\\
\textit{Escuela Nacional de Biblioteconomía y Archivonomía (ENBA – IPN)} — Ciudad de México
\begin{itemize}
  \item Planificación y adquisición de recursos materiales para sesiones prácticas;
        coordinación con proveedores de equipo tecnológico.
  \item Diseño de contenido académico sobre innovación tecnológica aplicada a procesos
        organizacionales y gestión de información.
\end{itemize}

\vspace{4pt}
% --- Cargo 4 ---
\noindent
\textbf{Desarrollador de Aplicaciones Móviles y VR/AR} \hfill
\textcolor{gris}{2023 – 2024}\\
\textit{Proyecto Académico / Profesional}
\begin{itemize}
  \item Evaluación y selección de herramientas y licencias de software; comparativa de
        costos y funcionalidades entre proveedores (Unity, Unreal, servicios cloud).
  \item Demostración tecnológica en TV Azteca México (2024): coordinación logística de
        equipos, materiales y proveedores externos para la producción.
\end{itemize}

% ═══════════════════════════════════════════════════════════════════════════════
%  CONOCIMIENTOS TÉCNICOS
% ═══════════════════════════════════════════════════════════════════════════════
\section{Conocimientos Técnicos}

\begin{tabularx}{\linewidth}{@{} >{\bfseries}p{5cm} X @{}}
  Mecánica y Manufactura &
    Lectura de planos, tolerancias, metrología dimensional, análisis estructural,
    análisis térmico, FEA (ANSYS), SolidWorks, MATLAB/Simulink \\[4pt]

  Gestión de Insumos y Materiales &
    Especificaciones técnicas de componentes mecánicos/electrónicos, comparativa de
    cotizaciones, control de calidad de recepción, cadena de suministro \\[4pt]

  Instrumentación y Sensores &
    Sensores de presión, temperatura y flujo; DAQ (LabVIEW/Python);
    control PID; normativas de instrumentación industrial \\[4pt]

  Análisis de Datos &
    Python (pandas, numpy, matplotlib), MATLAB, Excel/hojas de cálculo \\[4pt]

  Herramientas de Gestión &
    Órdenes de compra, evaluación de proveedores, logística de entrega,
    cumplimiento de políticas y código de ética empresarial \\[4pt]

  Idiomas &
    Español (nativo) · Inglés (intermedio — lectura técnica fluida) \\
\end{tabularx}

% ═══════════════════════════════════════════════════════════════════════════════
%  PROYECTOS Y LOGROS DESTACADOS
% ═══════════════════════════════════════════════════════════════════════════════
\section{Proyectos y Logros Destacados}

\begin{itemize}
  \item \textbf{Publicación IAC 2024 — Small Satellite Missions:}
        Coautor en el 75th International Astronautical Congress, Milán, Italia
        (DOI: 10.52202/078365-0120). Investigación en sistemas embebidos e instrumentación
        para cohetes de combustible sólido.
  \item \textbf{Banco de Pruebas para Cohetes de Combustible Sólido:}
        Caracterización, diseño y simulación de banco de pruebas mecánicas para medición
        de parámetros operativos. Publicado en \textit{Revista Hacia el Espacio} (2024).
        Involucró adquisición y validación técnica de sensores y materiales especializados.
  \item \textbf{Demostración VR en TV Azteca México (2024):}
        Coordinación de recursos técnicos y logísticos para demostración nacional de
        tecnología de Realidad Virtual aplicada a capacitación industrial.
\end{itemize}

\end{document}
